\Askhsh{34502}{2}
\begin{ylh}{1}
    \item 3.1. Είδη και στοιχεία τριγώνων
    \item 3.2. 1ο Κριτήριο ισότητας τριγώνων
    \item 3.12. Τριγωνική ανισότητα.
\end{ylh}
% ===================================================================
%  Αρχείο: 34502.tex (Fragment)
%  Αυτόματη μετατροπή μέσω doc2latex.py
% ===================================================================

ΘΕΜΑ 2

Στο ακόλουθο σχήμα, η ΑΔ είναι διάμεσος του τριγώνου ΑΒΓ και το Ε είναι σημείο στην προέκταση της ΑΔ, ώστε ΔΕ=ΑΔ.

Να αποδείξετε ότι:

\begin{alist}
\item ΑΒ = ΓΕ, \monades{12}

\item ΑΕ \textless{} ΑΒ + ΑΓ. \monades{13}


\begin{center}
\begin{tikzpicture}[scale=1]
\tkzDefPoint(0,0){B}
\tkzDefPoint(5,0){C}
\tkzDefPoint(1,3){A}
\tkzDefMidPoint(B,C) \tkzGetPoint{D}
\tkzDefPointBy[homothety=center D ratio -1](A) \tkzGetPoint{E}

\draw[l3] (A)--(B)--(C)--cycle;
\draw[l3] (A)--(E);
\draw[l3] (B)--(E);
\draw[l3] (C)--(E);

\tkzDrawPoints(A,B,C,D,E)
\tkzLabelPoint[above](A){$A$}
\tkzLabelPoint[left](B){$B$}
\tkzLabelPoint[right](C){$\Gamma$}
\tkzLabelPoint[below right](D){$\Delta$}
\tkzLabelPoint[below](E){$E$}

\tkzMarkSegments[mark=s|](B,D D,C)
\tkzMarkSegments[mark=s||](A,D D,E)
\end{tikzpicture}
\end{center}
\end{alist}
